%! Author = ben
%! Date = 27.11.2023

% Preamble
\documentclass[../main.tex]{subfiles}

% Packages

% Document
\begin{document}
    \chapter{Fazit und Ausblick}
    Mit SaveWorld haben wir es geschafft, eine moderne Plattform zu schaffen, die es jungen Menschen ermöglicht, sich mit dem Thema Nachhaltigkeit in vielen Facetten auseinanderzusetzen.
    Neben der reinen Wissensvermittlung, ist es uns ebenfalls gelungen, verschiedene Werkzeuge zu entwickeln, die den Benutzer in die Lage versetzt, ihr eigenes Leben nachhaltiger zu gestalten.
    Hierzu zählt beispielsweise der Öko-Tracker, welcher es ermöglicht, umweltschädliche Aktionen einzutragen und individuelle Ziele zur Verbesserung zu erhalten.
    \\
    Neben der Verbesserung des eigenen Lebens war es uns auch sehr wichtig, eine soziale Komponente einzubauen.
    Diese ist vor allem durch die integrierte Planung von ökologischen Projekten mit anderen Benutzern zusammen gegeben.
    Zudem ermöglicht das Forum die direkte Kommunikation zwischen Benutzern, sodass eventuelle Fragen beantwortet werden können.
    Durch die Möglichkeit, Chat-Gruppen zu erstellen, können Benutzer gemeinsam neue Themengebiete erkunden und sich gegenseitig motivieren.
    \\
    Mit den Artikeln ist es uns ebenfalls gelungen, ein breites Spektrum an Themen abzudecken und somit ein möglichst umfassendes Bild zu vermitteln, das auch ständig erweitert und aktualisiert wird.
    \\
    SaveWorld ist ein Projekt, welches sich stetig weiterentwickeln wird.
    In der Zukunft werden wir unseren Fokus auf folgende Aspekte legen:
    \begin{itemize}
        \item Wir werden die Oberfläche der App noch intuitiver gestalten, um die Benutzerfreundlichkeit zu erhöhen.
        \item Die Erweiterung der Inhalte hat oberste Priorität.
        Hierzu zählen sowohl die Artikel, als auch die Videos.
        Längere Videos im Querformat werden ebenfalls hinzugefügt, um einen tieferen Einblick in die Themen zu ermöglichen.
        \item Wettbewerbe und Challenges werden in der Zukunft ebenfalls hinzugefügt, sodass durch die Interaktion mit anderen zusätzliche Motivation entsteht.
        \item Eine Integration für Schulen ist ebenfalls geplant, sodass Lehrer die Möglichkeit haben, die App im Unterricht zu verwenden und zum Beispiel Klassen in Wettbewerben gegeneinander antreten zu lassen.
        \item Wir werden weitere Rechner (für die Berechnung von Emissionen) hinzufügen und die Möglichkeit, verschiedene Rechnerergebnisse mit einander zu vergleichen (z.B.\ Fahrt mit dem Auto vs.\ Fahrt mit dem Zug).
        \item Wir möchten kleine Spiele hinzufügen, um die App noch interaktiver zu gestalten und auch eine spielerische Komponente zu bieten.
        \item Die App wird ebenfalls in weitere Sprachen übersetzt, um möglichst viele Menschen zu erreichen.
        Hierbei ist es uns wichtig, dass die Übersetzungen nicht nur auf die Oberfläche beschränkt sind, sondern auch die Videos und Artikel übersetzt werden.
    \end{itemize}
\end{document}
